\documentclass[11pt]{article}
\usepackage[margin=1in]{geometry}
\usepackage{amsmath,amssymb,amsfonts,mathtools}
\usepackage{array,booktabs,enumitem,graphicx,xcolor,hyperref}
\usepackage[T1]{fontenc}
\usepackage[utf8]{inputenc}
\title{Effective Superstrings}
\author{Source-backed arXiv sample}
\date{}
\begin{document}
\maketitle
\begin{abstract}
We generalize the method of quantizing effective strings proposed by Polchinski and Strominger to superstrings. The Ramond-Neveu-Schwarz string is different from the Green-Schwarz string in non-critical dimensions. Both are anomaly-free and Poincare invariant. Some implications of the results are discussed. The formal analogy with 4D (super)gravity is pointed out.
\end{abstract}
\section{References And Citations}
\label{sec:refs}
This sample cites source keys \cite{pol} and \cite{div}. Section~\ref{sec:refs} and Equation~\ref{eq:source-demo} provide cross-reference coverage.
\begin{equation}
\label{eq:source-demo}
a^2 + b^2 = c^2
\end{equation}
\begin{thebibliography}{9}
\bibitem{pol} Source bibliography entry 1 extracted from arXiv metadata/source key `pol`.
\bibitem{div} Source bibliography entry 2 extracted from arXiv metadata/source key `div`.
\bibitem{hughes} Source bibliography entry 3 extracted from arXiv metadata/source key `hughes`.
\bibitem{gs} Source bibliography entry 4 extracted from arXiv metadata/source key `gs`.
\end{thebibliography}

\end{document}
